\documentclass[a4paper,12pt]{article}

\usepackage[utf8]{inputenc}
\usepackage{amsmath}
\usepackage{amssymb}
\usepackage{geometry}

\geometry{margin=2.5cm}

\title{IoT Homework -- Exercise 1}
\author{Angelica Ortiz 11191225, Daniel Nägele 11180610}
\date{May 2026}

\begin{document}

\maketitle
\thispagestyle{empty}
\newpage

\tableofcontents
\newpage

\section{Probability Mass Function}

The number of vibration peaks detected in each 10-second window follows a Poisson distribution with parameter:

\[
\lambda = 0.4
\]

The probability mass function of a Poisson random variable is:

\[
P(k)=\frac{\lambda^k e^{-\lambda}}{k!}
\]

Therefore:

\[
P(r=r_0)=P(k=0)=e^{-0.4}=0.6703
\]

\[
P(r=r_1)=P(k=1)=0.4e^{-0.4}=0.2681
\]

\[
P(r=r_2)=P(k>1)
\]

\[
P(r=r_2)=1-P(k=0)-P(k=1)
\]

\[
P(r=r_2)=1-0.6703-0.2681=0.0616
\]

Thus, the PMF is:

\[
P(r=r_0)=0.6703
\]

\[
P(r=r_1)=0.2681
\]

\[
P(r=r_2)=0.0616
\]

\section{Output Rates}

The reporting period is:

\[
T_{report}=10s
\]

If no vibration peaks are detected, the payload size is:

\[
512B
\]

Therefore:

\[
r_0=\frac{512\cdot8}{10}
\]

\[
r_0=409.6\ bit/s
\]

If one vibration peak is detected, the payload size is:

\[
2KB=2000B
\]

Thus:

\[
r_1=\frac{2000\cdot8}{10}
\]

\[
r_1=1600\ bit/s
\]

If more than one vibration peak is detected, the payload size is:

\[
4KB=4000B
\]

Thus:

\[
r_2=\frac{4000\cdot8}{10}
\]

\[
r_2=3200\ bit/s
\]

\section{Beacon Interval}

According to the standard IEEE 802.15.4 formulation, the Beacon Interval is computed from the minimum required rate.

The minimum output rate is:

\[
r_{min}=409.6\ bit/s
\]

Given:

\[
L=128B
\]

the Beacon Interval is:

\[
BI=\frac{L\cdot8}{r_{min}}
\]

Therefore:

\[
BI=\frac{128\cdot8}{409.6}
\]

\[
BI=2.5s
\]

This means that the reporting period of 10 seconds is divided into:

\[
\frac{10}{2.5}=4
\]

Beacon Intervals.

\section{CFP Slot Assignment}

The system operates with IEEE 802.15.4 in beacon-enabled mode using CFP only.

Given:

\[
R=250kbps
\]

\[
L=128B
\]

The slot duration is:

\[
T_s=\frac{L\cdot8}{R}
\]

\[
T_s=\frac{128\cdot8}{250000}
\]

\[
T_s=0.004096s
\]

\[
T_s=4.096ms
\]

The CFP is dimensioned according to the worst-case traffic scenario.

The largest payload generated by one node is:

\[
4000B
\]

The number of packets required over the full reporting period is:

\[
N_{packets}=\frac{4000}{128}
\]

\[
N_{packets}=31.25
\]

Approximating to the next integer:

\[
N_{packets}=32
\]

Since the reporting period spans 4 Beacon Intervals, the required number of slots per BI is:

\[
N_{slots,node}=\frac{32}{4}
\]

\[
N_{slots,node}=8
\]

The monitoring system contains 3 vibration monitoring nodes.

Therefore:

\[
N_{CFP}=3\cdot8
\]

\[
N_{CFP}=24
\]

Thus, the CFP slot assignment is:

\[
8\ slots/node
\]

and the total number of CFP slots is:

\[
N_{CFP}=24
\]

\section{Active and Inactive Time}

The active time includes:

\begin{itemize}
    \item one beacon slot;
    \item the CFP slots.
\end{itemize}

Therefore:

\[
N_{active}=24+1=25
\]

The active time is:

\[
T_{active}=25\cdot4.096ms
\]

\[
T_{active}=102.4ms
\]

The inactive time is:

\[
T_{inactive}=BI-T_{active}
\]

\[
T_{inactive}=2500ms-102.4ms
\]

\[
T_{inactive}=2397.6ms
\]

\[
T_{inactive}=2.3976s
\]

\section{Duty Cycle}

The duty cycle is computed as:

\[
DC=\frac{T_{active}}{BI}
\]

\[
DC=\frac{102.4}{2500}
\]

\[
DC=0.04096
\]

\[
DC=4.096\%
\]

\section{Additional Nodes}

The maximum allowed duty cycle is:

\[
DC_{max}=10\%
\]

The active time is:

\[
T_{active}=(1+8M)T_s
\]

where \(M\) is the number of vibration monitoring nodes.

Imposing the duty-cycle constraint:

\[
\frac{(1+8M)\cdot4.096}{2500}<0.1
\]

\[
1+8M<61.035
\]

\[
8M<60.035
\]

\[
M<7.50
\]

Therefore, the maximum integer number of nodes is:

\[
M_{max}=7
\]

The system already contains 3 nodes.

Therefore:

\[
\Delta M=7-3
\]

\[
\Delta M=4
\]

Hence, 4 additional vibration monitoring nodes can be added while keeping the duty cycle below 10\%.

\[
\Delta M=4\ additional\ nodes
\]

\end{document}
